\hypertarget{appendix-e89a8f290ef5}{}
\section{Atomic publication and durability}\label{appendix-e89a8f290ef5-atomic-publication-and-durability}

\textbf{Curriculum depth:} Introduced · \textbf{First seen:} C003 · \textbf{Latest development:} C003

\subsection{Definition and mechanism}\label{appendix-e89a8f290ef5-definition-and-mechanism}

Publishing a prepared temporary file by replacement concerns observers of the filesystem namespace. Crash durability concerns whether contents and namespace updates survive failures; it needs a specified platform/filesystem protocol. Neither follows from building a vector of bytes. The supplied AtomicFile excerpt is an unverified repository claim, not a tested durability guarantee.

\subsection{Prerequisites}\label{appendix-e89a8f290ef5-prerequisites}

\begin{itemize}
\tightlist
\item
  \hyperref[appendix-91d09ebed979]{Serialization} --- Publication persists already-specified bytes, rather than defining their format.
\item
  \href{../05-Engine-I/main.pdf\#nameddest=appendix-d52e4486f527}{Ownership and non-owning views (05-Engine-I, appendix-d52e4486f527)} --- A mapping or publication pipeline must keep referenced bytes valid.
\end{itemize}

\subsection{Used by}\label{appendix-e89a8f290ef5-used-by}

Update downstream uses when later chunks build on this concept.

\subsection{Recall and next depth}\label{appendix-e89a8f290ef5-recall-and-next-depth}

Name the observers and failure model for an atomic guarantee.

\textbf{Pending:} Platform replacement semantics, flush ordering, recovery and cross-volume limits.

\subsection{Source and navigation}\label{appendix-e89a8f290ef5-source-and-navigation}

\hyperref[01-foundations-i-ch002-where-the-bytes-live]{C003 teaching chapter} · \hyperref[appendix-fd090d982a15-m003---parts-41-60]{Part definitions} · \href{../02-Foundations-II/main.pdf\#nameddest=appendix-4ad47bdd269a}{Dependency Map (02-Foundations-II, appendix-4ad47bdd269a)}

\subsection{Executed POSIX durability sequence}\label{appendix-e89a8f290ef5-executed-posix-durability-sequence}

The durability contract now has an actual syscall-level witness in \texttt{tools/labs/atomic\_publication\_durability\_lab.cpp}.

The reference creates a staging file with exclusive creation, writes the complete payload, calls \texttt{fsync} on the file, closes it, renames it within the same directory, opens the containing directory, and calls \texttt{fsync} on that directory before treating publication as complete. It then verifies the final bytes and absence of the staging name. Separate tests reject a pre-existing staging file and a traversal-style final name.

This is stronger than ``write temp then rename,'' because rename atomicity and durable persistence are different properties. The file sync requests persistence of file contents/metadata; the directory sync requests persistence of the directory-entry change. Successful calls on CI are evidence that the operating-system/filesystem accepted that ordering. They are \textbf{not} a laboratory power-cut test and cannot promise identical persistence guarantees on every filesystem, storage controller, Windows API, network filesystem or production asset pipeline.
